% !TITLE 一剪梅·红藕香残玉簟秋
% !AUTHOR 李清照
% !DYNASTY 两宋
% !GENRE 词
% !ORDER 14

\begin{poem}
红藕\textsuperscript{1}香残玉簟\textsuperscript{2}秋。轻解罗裳\textsuperscript{3}，独上兰舟。
云中谁寄锦书\textsuperscript{4}来？雁字\textsuperscript{5}回时，月满西楼。

花自飘零水自流。一种相思，两处闲愁\textsuperscript{6}。
此情无计可消除，才下眉头，却上心头\textsuperscript{7}。
\end{poem}

\begin{annotations}
\notetext{1}{红藕：红色的荷花。荷花香残，暗示秋天已到。}
\notetext{2}{玉簟：光滑如玉的竹席。簟（diàn），竹席。秋天竹席变凉，触觉上暗示季节转换。}
\notetext{3}{罗裳：丝质的裙子。"轻解罗裳"——轻轻脱下外裙（换上便装），独自登上小船。}
\notetext{4}{锦书：写在锦帛上的书信，指妻子给丈夫的信。用了前秦苏蕙织锦回文诗寄给丈夫窦滔的典故。}
\notetext{5}{雁字：大雁飞行的队列像"人"字或"一"字。大雁传书是古代的想象——看见大雁飞回来了，却不见书信，只有满楼的月光。}
\notetext{6}{两处闲愁：一种相思，牵动两处的闲愁——丈夫和妻子虽然分开，但思念是相互的。}
\notetext{7}{才下眉头，却上心头：刚从眉头上消除（不再皱眉），又涌上心头。愁不是表面表情能排解的——它埋在心里，任何时候都会冒出来。这两句是千古名句。}
\end{annotations}
